%-----------------------------------------------------------------------
% Threats to Validity subsection for Section VI (Discussion and Limitations)
% Insert after line 610 of main.tex, before \section{Related Work}
%-----------------------------------------------------------------------

\subsection{Threats to Validity}

We identify seven threats to the validity of the claims made in this paper and describe how the framework's design addresses each.

\textbf{Scanner Dependency (Construct Validity).}
Constraint Guard's category-specific rules require the upstream static analyzer to provide accurate vulnerability categorization. On all three real-world projects, clang-tidy assigned virtually all findings to the \textsc{unknown} category, preventing category-gated rules from firing and producing identical Precision@10 scores between the Baseline and Deterministic configurations. This is a real deployment limitation of the current evaluation setup. The LLM-assisted reclassification experiment addresses this threat directly: as reported in Section~\ref{sec:llm_reclassification}, the LLM recovered 403 of 463 \textsc{unknown} findings (87\%), producing 201 tier changes---84 findings from LOW to MEDIUM in Zephyr and 117 findings from MEDIUM to HIGH in esp-fc---and identifying one novel category, \texttt{logic\_error}, not present in the predefined taxonomy. When reclassification is enabled, category-specific rules fire on findings that a scanner-level miscategorization would otherwise silence.

\textbf{Beyond-Scanner Coverage (Construct Validity).}
Static analyzers are structurally limited in their ability to detect concurrency-specific embedded vulnerabilities such as race conditions, TOCTOU hazards, blocking calls in ISR context, and incorrect \texttt{volatile} usage. Findings of this type do not appear in the SARIF output and therefore cannot be scored by the deterministic engine. The LLM-surfaced candidate discovery experiment reported in Section~\ref{sec:llm_discovery} demonstrates a framework-level response: across the three real-world projects, the LLM identified 56 candidate findings, of which 52 were confirmed as new (not present in the existing SARIF output), yielding an overall precision of 92.9\%. These candidates are assigned scores and tiers through the same deterministic engine, extending coverage to vulnerability classes the scanner cannot reach.

\textbf{Ground Truth Scale (Internal Validity).}
The expert agreement evaluation in Section~\ref{sec:evaluation} is based on 30 manually labeled findings per project (90 total), labeled by the authors of this paper. This scale is insufficient for statistical generalization and introduces potential annotator bias in the labels themselves. The framework addresses this limitation through complementary validation: the controlled benchmark (Experiment 8) evaluates scoring behavior on eight deliberately introduced vulnerabilities with known ground truth across six CWE categories, providing a setting where the correct answer is unambiguous. The 70.9\% expert agreement on real-world projects, while modest in scale, establishes a baseline for practical triage alignment that can be replicated and extended by independent evaluators using the published evaluation harness.

\textbf{Evaluator Bias (Internal Validity).}
LLM enrichment quality was rated by the authors of this paper rather than independent evaluators, introducing potential confirmation bias into the usefulness scores (3.9/5 average) and evidence correctness measurements (80\% average). We mitigate this risk in two ways. First, evidence correctness is evaluated against a verifiable criterion: each citation in the LLM output references a specific source file and line range, and correctness is determined by checking whether that location exists and contains the described construct---a judgment that does not require domain expertise. Second, the structured output schema constrains LLM responses to factual claims about code (tags, fix suggestions, and cited line numbers), limiting the scope for subjective interpretation. Independent replication with external reviewers remains a priority for future work.

\textbf{LLM Overhead Exclusion (Internal Validity).}
The network-latency-bound cost of LLM enrichment (2--10 seconds per finding for GPT-5-mini) is excluded from the 4.02\,ms overhead measurement reported in Section~\ref{sec:evaluation}. This is a design decision, not an oversight. The deterministic scoring pipeline---constraint extraction, scoring, and reporting---constitutes the CI/CD gate: it runs on every commit and produces the pass/fail exit code used by pipeline policies. LLM enrichment is an optional post-gate layer configured separately via environment variables; it does not execute on the blocking path unless explicitly enabled. The 4.02\,ms figure therefore measures the latency of the decision that blocks or permits a build, which is the operationally relevant quantity for CI/CD deployment.

\textbf{Hardcoded Rule Thresholds (Construct Validity).}
The 14 scoring rules use fixed hardware thresholds---stack budget 4{,}096\,bytes, RAM 64\,KB, ISR latency 50\,$\mu$s and 100\,$\mu$s---that may not capture the full range of embedded platform constraints. These values represent common boundaries in safety-critical CPS design, but a platform operating outside these ranges might produce misleading scores. The controlled benchmark (Experiment 8) evaluates scoring behavior across four hardware profiles spanning from 512\,KB RAM to 20\,KB RAM and 500\,$\mu$s to 50\,$\mu$s ISR latency, demonstrating that the engine produces meaningfully differentiated priorities as constraints tighten. Rule thresholds are derived from the \texttt{HardwareSpec} fields provided in the user configuration; the rule-based architecture supports adding or modifying rules without changes to the scoring engine, enabling platform-specific calibration.

\textbf{Generalization (External Validity).}
The evaluation covers three CPS projects (FreeRTOS-Kernel, Zephyr RTOS, esp-fc) running on ARM Cortex-M3 and ESP32 hardware, analyzed with clang-tidy and the Clang Static Analyzer. The explicit rule-to-category mappings in \texttt{sarif\_rule\_map.py} cover 30+ Clang SA checker IDs; analyzers with different rule namespaces will produce a higher proportion of \textsc{unknown} findings under the current mappings. However, the SARIF ingestion layer is tool-agnostic: any analyzer producing standards-compliant SARIF output can be processed, with CWE-based category resolution as a secondary assignment path. The LLM reclassification layer further reduces dependence on scanner-specific mappings by recovering semantic categories from finding descriptions. Extending the evaluation to additional analyzers, hardware architectures, and application domains beyond CPS is identified as the primary direction for future work.

% HUMANIZED VERSION:

\subsection{Threats to Validity}

We list seven threats to validity of the claims made in this paper and how the framework design attempts to mitigate them.

\textbf{Scanner Dependency (Construct Validity).}
Constraint Guard's category-specific rules rely on the upstream static analyzer to categorize vulnerabilities accurately. On all three real-world projects, clang-tidy assigned almost all findings to the \textsc{unknown} category so that our category-gated rules could not fire. This produced identical scores for the Baseline and Deterministic configurations on Precision@10. This is a limitation of the real deployment of our current evaluation setup. The threat is mitigated by our LLM-assisted reclassification experiment, reported in Section~\ref{sec:llm_reclassification}: the LLM recovered 403 out of 463 \textsc{unknown} findings (87\%), producing 201 tier changes---84 findings from LOW to MEDIUM in Zephyr and 117 findings from MEDIUM to HIGH in esp-fc---and produce one novel category, \texttt{logic\_error}, that does not already exist in the project-defined taxonomy. With LLM reclassification enabled, our category-specific rules fire on findings that high-level scanner miscategorization would silence.

\textbf{Beyond-Scanner Coverage (Construct Validity).}
Static analyzers are to some degree limited from being able to detect concurrency-specific embedded vulnerabilities such as race conditions, TOCTOU hazards, blocking calls in ISR context, and incorrect \texttt{volatile} usage. Findings of this type not present in the SARIF file cannot be scored by the deterministic engine. The user-generated candidate discovery experiment reported in Section~\ref{sec:llm_discovery} shows the framework-level response: the LLM produced 56 candidate findings, of which 52 were confirmed as new (not present in existing SARIF output), producing a precision of 92.9\% over the remaining 34 candidate findings. These candidates are also given scores and tiers based on the same deterministic engine, mitigating the aspects of the vulnerability class that the scanner itself cannot reach.

\textbf{Ground Truth Scale (Internal Validity).}
As discussed in Section~\ref{sec:evaluation}, expert agreement is tested based on 30 manually labeled findings per project (90 total), by the authors of this paper. This scale fails to hit the threshold needed for statistical generalization and subjects the labels themselves to possible annotator bias. Rather than attempting generalization based on this limited scale, the framework also incorporates reassurance through complementary validation results: the controlled benchmark (Experiment 8) evaluates scoring behavior on eight intentionally introduced vulnerabilities with known ground truth across six CWE categories, in the narrow design space where the right answer is clear. The 70.9\% expert agreement on real projects, admittedly modest in scale, does offer a baseline to which others can compare practical triage alignment. Others can reproduce these psychometric results using our published evaluation harness.

\textbf{Evaluator Bias (Internal Validity).}
We rated, ourselves, the LLM enrichment quality rather than using independent evaluators. This introduces possible confirmation bias into the usefulness scores reported (3.9/5 average) and the measurements of evidence correctness (80\% average) in Section~\ref{sec:evaluation}. We ameliorate this risk in two ways. First, we evaluate evidence correctness against a verifiable criterion: each citation in our LLM output cites a particular source file and range of lines, and evidence correctness is also based on whether one can find a file in that location, and therein locate the described construct---the result of an absolute judgment not requiring subject matter expertise. Second, the structured output schema constrains LLM responses to factual claims about code (tags, fix suggestions, and cited line numbers), limiting the scope for subjective interpretation. Independent replication with external reviewers remains a priority for future work.

\textbf{LLM Overhead Exclusion (Internal Validity).}
The network-latency-bound cost of LLM enrichment (2--10 seconds per finding for GPT-5-mini) is excluded from the 4.02\,ms overhead measurement reported in Section~\ref{sec:evaluation}. This is a design decision, not an oversight. The deterministic scoring pipeline---constraint extraction, scoring, and reporting---constitutes the CI/CD gate: it runs on every commit and produces the pass/fail exit code used by pipeline policies. LLM enrichment is an optional post-gate layer configured separately via environment variables; it does not execute on the blocking path unless explicitly enabled. The 4.02\,ms figure therefore measures the latency of the decision that blocks or permits a build, which is the operationally relevant quantity for CI/CD deployment.

\textbf{Hardcoded Rule Thresholds (Construct Validity).}
The 14 scoring rules use fixed hardware thresholds---stack budget 4{,}096\,bytes, RAM 64\,KB, ISR latency 50\,$\mu$s and 100\,$\mu$s---that may not capture the full range of embedded platform constraints. These values represent common boundaries in safety-critical CPS design, but a platform operating outside these ranges might produce misleading scores. The controlled benchmark (Experiment 8) evaluates scoring behavior across four hardware profiles spanning from 512\,KB RAM to 20\,KB RAM and 500\,$\mu$s to 50\,$\mu$s ISR latency, demonstrating that the engine produces meaningfully differentiated priorities as constraints tighten. Rule thresholds are derived from the \texttt{HardwareSpec} fields provided in the user configuration; the rule-based architecture supports adding or modifying rules without changes to the scoring engine, enabling platform-specific calibration.

\textbf{Generalization (External Validity).}
The evaluation covers three CPS projects (FreeRTOS-Kernel, Zephyr RTOS, esp-fc) running on ARM Cortex-M3 and ESP32 hardware, analyzed with clang-tidy and the Clang Static Analyzer. The explicit rule-to-category mappings in \texttt{sarif\_rule\_map.py} cover 30+ Clang SA checker IDs; analyzers with different rule namespaces will produce a higher proportion of \textsc{unknown} findings under the current mappings. However, the SARIF ingestion layer is tool-agnostic: any analyzer producing standards-compliant SARIF output can be processed, with CWE-based category resolution as a secondary assignment path. The LLM reclassification layer further reduces dependence on scanner-specific mappings by recovering semantic categories from finding descriptions. Extending the evaluation to additional analyzers, hardware architectures, and application domains beyond CPS is identified as the primary direction for future work.